\section*{第 \ref{ch:b2:enolates-aldol} 章——烯醇、烯醇负离子与羟醛反应}

\begin{solution}{exo:b2:enolates-aldol:1}
丁酮：丁-1-烯-2-醇 \ce{CH2=C(OH)CH2CH3} 与丁-2-烯-2-醇 \ce{CH3C(OH)=CHCH3}（(\textit{E}) 与
(\textit{Z})）。环己酮：环己-1-烯-1-醇。
\end{solution}

\begin{solution}{exo:b2:enolates-aldol:2}
乙烷 $<$ 丙酮 $<$ 3-氧代丁酸乙酯（10.7）$<$ 戊烷-2,4-二酮（9.0）。
\end{solution}

\begin{solution}{exo:b2:enolates-aldol:3}
3-羟基丁醛，加热后得丁-2-烯醛。
\end{solution}

\begin{solution}{exo:b2:enolates-aldol:4}
丙二酸二乙酯、乙醇钠、1-溴丁烷；水解后得丁基丙二酸；
加热脱去 \ce{CO2}：\ce{CH3(CH2)3CH2COOH}，己酸。
\end{solution}

\begin{solution}{exo:b2:enolates-aldol:5}
动力学烯醇负离子：在 C6 处（\ce{CH2} 一侧）去质子，使用 LDA，于 \qty{-78}{\degreeCelsius} 下进行。热力学烯醇负离子：
生成取代程度更高的双键 C1=C2（朝向带甲基的碳），
以醇盐为碱，在其醇溶剂中于室温下进行。
\end{solution}

\begin{solution}{exo:b2:enolates-aldol:6}
4-苯基丁-3-烯-2-酮（亚苄基丙酮），\ce{C6H5CH=CHCOCH3}。苯甲醛不含 $\alpha$ 氢，它
不能形成烯醇负离子，只能作亲电试剂。
\end{solution}

\begin{solution}{exo:b2:enolates-aldol:7}
一个甲基（C1）形成的烯醇负离子进攻另一羰基（C5），生成五元环，
脱水后得 3-甲基环戊-2-烯-1-酮。另一种可能是 C3 处的烯醇负离子进攻 C5，
将生成张力很大的三元环。
\end{solution}

\begin{solution}{exo:b2:enolates-aldol:8}
\ce{2CH3COOC2H5 -> CH3COCH2COOC2H5 + C2H5OH}，在乙醇钠作用下进行；产物被去质子化，
酸化后得到 3-氧代丁酸乙酯。
\end{solution}

\begin{solution}{exo:b2:enolates-aldol:9}
丙酮、乙醛和苯乙酮（它们含有 \ce{CH3-CO} 基团；乙醛为 \ce{CH3CHO}）。
戊-3-酮则没有。
\end{solution}

\begin{solution}{exo:b2:enolates-aldol:10}
一端的烯醇负离子进攻另一端的酯基，生成 2-氧代环戊烷-1-羧酸乙酯，
一个五元环。
\end{solution}

\begin{solution}{exo:b2:enolates-aldol:11}
在 3-甲基戊-2-酮中，手性碳原子（C3）是 $\alpha$ 碳：烯醇化使其变为平面结构，
两个面上的再质子化都会使其外消旋化。在 4-甲基己-2-酮中，手性碳原子位于 C4，
位于羰基的 $\beta$ 位，不受烯醇化影响。
\end{solution}

\begin{solution}{exo:b2:enolates-aldol:12}
丙酮与两当量苯甲醛在碱中反应。若要用两种不同的醛，应
先让丙酮与一种醛缩合（取一当量，分离得亚苄基丙酮），再
使其剩余甲基与第二种醛缩合。若一锅投料，两种醛相互竞争，
会得到三种二烯酮的混合物。
\end{solution}

\begin{solution}{pb:b2:enolates-aldol:1}
\textbf{1.} 丙酮，含六个 $\alpha$ 氢（每个甲基上三个）。
\textbf{2.} $\mathrm{CH_3{-}CO{-}CH_2^-} \leftrightarrow \mathrm{CH_3{-}C(O^-){=}CH_2}$.
\textbf{3.} 与羰基相邻的碳上没有氢（该碳属于环且
不带氢）。
\textbf{4.} 过量的苯甲醛是更好的亲电试剂，丙酮的烯醇负离子
与它相遇的机会多于与丙酮羰基相遇。
\textbf{5.} 酮有两个碳基团将电子云推向羰基碳，
阻碍试剂接近；醛只有一个。
\textbf{6.} 只需少量烯醇负离子：反应时它会重新生成。
\textbf{7.} 烯醇负离子的碳与苯甲醛的羰基碳成键，生成
4-羟基-4-苯基丁-2-酮的醇盐。
\textbf{8.} 脱去一个 $\alpha$ 氢，并从 $\beta$ 碳上失去氢氧根：4-苯基丁-3-烯-2-酮。
\textbf{9.} 所形成的双键与苯环和羰基共轭。
\textbf{10.} 另一个甲基仍有三个 $\alpha$ 氢。
\textbf{11.} \ce{2C6H5CHO + CH3COCH3 -> C6H5CH=CHCOCH=CHC6H5 + 2H2O}.
\textbf{12.} 脱水反应，以及产物沉淀而离开溶液。
\textbf{13.} 两种，均为 (\textit{E})。
\textbf{14.} (\textit{Z}) 型同分异构体使苯基靠近羰基，产生空间张力。
\textbf{15.} 两个 \ce{C=C}、一个 \ce{C=O} 以及两个苯环：贯穿整个分子的共轭体系。
\textbf{16.} 其扩展的共轭体系缩小了 $\pi$ 能隙（见\cref{prop:b2:huckel:gap}）：它在
可见光的蓝端吸收，看起来呈黄色。
\textbf{17.} 否：它为平面结构（或近似平面），具有对称面。
\textbf{18.} 106.12、58.08 与 \qty{234.30}{g/mol}。
\textbf{19.} 苯甲醛 $2.1 \times 1.045/106.12 = \qty{20.7}{mmol}$；丙酮 $0.75 \times
0.7845/58.08 = \qty{10.1}{mmol}$。
\textbf{20.} 丙酮需要 \qty{20.3}{mmol} 苯甲醛，而可用量为 \qty{20.7}{mmol}：
丙酮为限量试剂，仅勉强足量。
\textbf{21.} $0.01013 \times 234.30 = \qty{2.37}{g}$.
\textbf{22.} 剩余的苯甲醛留在乙醇中，被洗去。
\textbf{23.} 单缩合产物，即亚苄基丙酮。
\textbf{24.} 二亚苄基丙酮的量为 $0.75 \times 2.37 = \boldsymbol{\approx \qty{1.78}{g}}$。
\end{solution}
